\section{评估设置}
\label{sec:setup}

为评估 \hesp{}，我们围绕四个问题组织评估：
\begin{enumerate}
  \item \emph{可靠性：}\hesp{} 能否让小型本地模型成为可靠的分诊智能体？它是否依赖对环境的神谕知识？（\S\ref{sec:rq-help}）
  \item \emph{排序还是剥夺选择：}收益来自信息增益排序，还是仅仅来自把选探针的权力从模型手中拿走？（\S\ref{sec:rq-rank}）
  \item \emph{探测还是停止：}“选择探什么”和“决定何时停”是否是两种独立的失败？每个模型各有哪一种？（\S\ref{sec:rq-stop}）
  \item \emph{失败方式：}从安全角度看，结论是如何出错的？（\S\ref{sec:security}）
\end{enumerate}
每个问题都对应一项预注册研究，每项研究都在上一项运行完成之后才设计（\cref{tab:studies}）。\S\ref{sec:realdata} 进一步追问：公开安全数据能否检验同样的问题。

\begin{table*}[t]
  \caption{四项在 sec-triage 上的预注册研究总览，共 7{,}272 个回合。仓库中按版本号标记这些研究。}
  \label{tab:studies}
  \centering\small
  \begin{adjustbox}{max width=\textwidth}
  \begin{tabular}{llllrl}
    \toprule
    研究 & 版本 & 问题 & 模型 & 回合数 & 主要终点 \\
    \midrule
    初始研究 & v0.5 & 控制器有没有用？ & 3 个 Qwen & 1{,}944 & 完整 \hesp{} 对 Memory-only，逐模型 \\
    表来源研究 & v0.6 & 表必须来自神谕吗？ & 3 个 Qwen & 1{,}728 & 计数表 \hesp{} 对 Memory-only，7B \\
    分解研究 & v0.8 & 排序还是剥夺选择？加入第二个家族 & 5 & 1{,}800 & 排序对随机（Q-7B）；\hesp{} 对 Memory-only（L-8B） \\
    停止研究 & v0.9 & 探什么 vs.\ 何时停 & 5 & 1{,}800 & 控制器停止；给定停止后的选择（均为 L-8B） \\
    \bottomrule
  \end{tabular}
  \end{adjustbox}
\end{table*}

\subsection{环境与表}

\subsubsection{sec-triage 环境} 每个回合是一个 Web 服务上的一条告警。隐藏起因是八种解释之一：三种攻击（撞库、SQL 注入探测、DNS 命令与控制）、一名正在外泄数据的内部人员、一个暴露的管理面板、一个已知有漏洞的组件、一次经授权的漏洞扫描，以及一次监控误报。假设集合再加上\textit{其他}。十个只读探针，每个成本 1 到 3 单位，覆盖认证失败、HTTP 模式、来源信誉、管理面板暴露、按用户的出站流量、DNS 特征、组件清单、变更工单、威胁情报和一本通用操作手册。我们有意让捷径失效。威胁情报只会说某个指标是恶意的，而三种攻击共享这个答案。操作手册不含任何信息。有几种良性起因在运行正确的探针之前看起来都像攻击。每种起因有三个变体。\textit{base} 中，结果是确定的，只有 10\% 的威胁情报延迟。\textit{noise} 增加 15\% 的暂时性 HTTP 503 故障和 30\% 的延迟。\textit{drift} 中，事件在第二个探针之后被替换：状态版本递增，真实起因改变。由此得到 24 个任务，每项研究在每种配置下把每个任务运行三次。每个回合由它自己的回环 HTTP 沙箱提供服务。

\subsubsection{预测表} 我们比较 $P(o \mid h, a)$ 的三种来源。\emph{设计者}表来自环境的生成函数，作为神谕参照。\emph{计数}表（表格中记为 “empirical”）按 \S\ref{sec:tables} 的方法估计，每个起因和变体用 $k \in \{1, 5, 20, 100\}$ 个开发回合，只用 base 和 noise 变体，种子取自与评估不相交的空间。\emph{模型自报}表由规划器 LLM 自己给出。每张表都在任何评估回合运行之前冻结并计算哈希。由于开发回合和评估回合来自同一个生成器，计数表告诉我们需要多少观测才够，而不能说明方法在陌生环境中的表现。

\begin{figure*}[t]
  \centering
  \includegraphics[width=\textwidth]{fig_motivating.pdf}
  \caption{同一条告警、同一个本地模型、四种划分调查的方式。告警的真实起因是暴露的管理面板；每一行的规划器都是 Llama-3.1-8B，种子相同。每个小块是一个实际运行的探针（虚线：被控制器拦截的重复提议）；红色标出确认起因的那条观测。当由模型决定何时停止时（前两行），它拿到了决定性证据，却继续探测直到预算耗尽。控制器侧停止会下结论（后两行），而信息增益选择只用三个探针就做到了。四行均原样取自停止研究的归档日志。}
  \label{fig:motivating}
\end{figure*}

\subsection{实验设计}

\subsubsection{配置} 所有配置使用相同的模型、提示模板、探针目录、预算（10 次工具调用、10 个成本单位、12 次决策、900\,秒）、重复拦截和验证器。\emph{ReAct-style}（ReAct 式）：规划器只看到原始观测历史，每个探针都由它选。\emph{Memory-only}（仅记忆）：规划器还能看到账本，但仍由它选每个探针。\emph{\hesp{}}：由控制器选探针（EIG/成本或随机），可展示或不展示排序，可带或不带结束守卫和控制器侧停止。\Cref{fig:motivating} 在同一条告警上展示了其中四种配置。停止研究之后，我们还用一个按目录顺序提议探针、从不结束的脚本替换 LLM，运行了冻结的控制器。这个\emph{无 LLM 参照}使用停止研究的种子、任务、预算和表（216 个回合）；它没有预注册，我们把它作为探索性结果报告。

\subsubsection{实现} \hesp{} 约 2{,}900 行 Python，只用标准库。规划器可插拔：同一个 JSON 决策协议同时服务 vLLM 和 Ollama 端点，以及无 LLM 运行所用的脚本化策略。每个回合有自己的回环 HTTP 沙箱，因此探针是真实的 HTTP 请求，带有暂时性故障和状态变化，任何回合都无法观察到另一个回合。每项研究都使用 Qwen2.5-7B-Instruct（bf16）、Qwen2.5-32B-Instruct-AWQ 和 Qwen2.5-72B-Instruct-AWQ；分解研究和停止研究另加 Llama-3.1-8B-Instruct（bf16）和 Llama-3.1-70B-Instruct-AWQ-INT4。模型由 vLLM 在单块 RTX~6000 级 GPU 上提供服务，温度 0.2，每次调用使用不同的种子。一个成对、打乱顺序、可续跑的运行器最多同时执行 32 个回合；如果清单或控制器源码的哈希发生变化，它会拒绝续跑。每个模型跑完后，它审计全部回合、归档日志，并记录归档文件的 SHA-256。停止研究的 1{,}800 个回合用时 28 分钟。一套 145 个单元测试覆盖账本、选择器、盲提示、守卫、控制器停止、审计，以及表估计器与环境生成器之间的隔离。

\subsubsection{评估方法} 每项研究的假设、配置、主要终点和分析方法，连同源码和表的哈希，都在其任何回合运行之前写下并冻结；冻结之后的改动以勘误形式追加（附录~\ref{app:prereg}）。结果指标是\emph{已验证完成率}，即以已验证结论结束的回合所占比例。其他任何结束方式，包括超时和格式错误的输出，都算失败。我们报告按任务配对的差值，以及百分位数的任务簇自助法区间（2{,}000 次重抽样，24 个簇）。当一项研究有两个主要终点时，每个终点使用 Bonferroni 校正后的 97.5\% 区间；所有次要比较都是探索性的，未做校正。我们还报告描述性的安全指标。有两种起因是良性的（经授权的扫描和监控误报），其余六种需要处置。在以结论结束的回合中，\emph{漏报攻击率}是真实起因需处置、而结论给出良性起因的比例；\emph{误升级率}则相反；\emph{错因率}是结论给错起因的比例。\emph{未解决率}是所有回合中没有已验证结论的比例。真值是给出结论那一刻生效的起因。所有表格和图都由脚本从原始结果文件生成。
